\section{Limitations}
\label{sec:limitations}

Intentionality is fixed at a single (knowing-but-non-malicious) level in
the core confirmatory design, so this study cannot yet speak to whether any
observed gender effect on fault attribution varies with the norm-violator's
stated intent; an intentionality-robustness arm is designed but deferred.
Relationship context (type and duration) is fixed per family rather than
crossed, so family-level effects cannot be cleanly separated from
relationship-stage differences that happen to co-occur with a given family
in this design. The severity manipulation has not yet been piloted against
the full 288-vignette core set -- only a 20-vignette internal review sample
has been read through by the research team -- so severity levels are
currently validated by the design team's judgment rather than by an
independent rater or manipulation check. The obligation-source taxonomy and
severity anchors were authored by the research team rather than derived
from or validated against a corpus of real disputes, and the four-gender
configurations (male/female agent $\times$ male/female partner) operationalize
gender as a strict binary, which does not capture non-binary or genderqueer
identities and should not be over-generalized beyond that binary
operationalization. The evaluated models are commercial and open-weight
systems accessed through a single API aggregator (OpenRouter) at a fixed
point in time; findings may not generalize to other models, model
versions, or direct-vendor deployments with different system prompts or
safety filtering. Because r/AmItheAsshole is heavily represented in
pretraining data, there is a standing risk that a model's fault-rating
output reflects memorized genre patterns rather than independent moral
reasoning \citep{russo2025pluralistic}; writing every vignette from
scratch in a fixed template (Section~\ref{sec:dataset}) avoids the
verbatim-memorization risk Russo et al.\ had to engineer around by
rewriting real posts, but the milder risk of \emph{structural} genre
familiarity -- recognizing the AITA-style framing and conventions even in
an invented scenario -- remains unaddressed, since memorization is known
to scale with model capacity and exposure frequency
\citep{carlini2023memorization} and is not fully caught by surface-level
contamination checks \citep{xu2024contamination}. A model-based novelty
check (prompting a separate call to guess whether a vignette is
real-vs-synthetic) is designed but not yet run.
